% sections/empirical_foundations.tex --- Empirical Foundations
\section{Empirical Foundations}
\label{sec:empirical}

\bench{} is calibrated against an empirical study of journal-entry
structure conducted on a large corpus of real general
ledgers~\citep{ivertowski2024apn}.  This section summarises the
findings that drive the parameters of the \system{} generator and the
distributional checks reported in Section~\ref{sec:evaluation}.  All
figures in this section are reproduced from
\citet{ivertowski2024apn} with permission of the author; the
underlying client data are confidential and are not redistributed.

\subsection{Corpus}
\label{sec:empirical:corpus}

The corpus comprises $155$ general-ledger datasets (approximately
$0.2\%$ of the relevant population), drawn from a representative
mix of industry sectors and covering each subject's full fiscal year
of activity.  Aggregate volumes are approximately $364$ million
journal entries and $2.4$ billion line items.  Datasets are
classified into three size categories
(Table~\ref{tab:emp:size}): small ($<\!1$M lines per annum,
$74$ sets), medium ($1$--$10$M lines per annum, $47$ sets), and
large ($>\!10$M lines per annum, $34$ sets).

\begin{table}[H]
  \centering
  \caption{Classification of the source datasets by size.  Source:
    \citet{ivertowski2024apn}.}
  \label{tab:emp:size}
  \begin{tabular}{lccc}
    \toprule
                  & \textbf{Small} & \textbf{Medium} & \textbf{Large}\\
                  & ($<\!1$M lpa)  & ($1$--$10$M lpa) & ($>\!10$M lpa)\\
    \midrule
    Number of sets & $74$ & $47$ & $34$\\
    \bottomrule
  \end{tabular}
\end{table}

\subsection{Line-item distribution per journal entry}
\label{sec:empirical:line-items}

The empirical distribution of line items per journal entry is
extremely concentrated.  $92.60\%$ of all postings carry between
two and nine line items; the remaining mass tapers across orders of
magnitude up to a maximum of $47\,059$ line items in a single
posting.  Table~\ref{tab:emp:bins} reproduces the bin-level
distribution.  Within bin~1 ($2$--$9$ line items),
Table~\ref{tab:emp:bin1} shows that two-line entries alone account
for $60.68\%$ of the corpus, with even-count entries
($2,4,6,\ldots$) cumulatively dominating: roughly $88\%$ of
journal entries carry an even number of lines.

\begin{table}[H]
  \centering
  \caption{Distribution of journal-entry line counts across the
    corpus.  Source: \citet{ivertowski2024apn}, Table~II.}
  \label{tab:emp:bins}
  \begin{tabular}{ccrr}
    \toprule
    \textbf{Bin} & \textbf{Line items} & \textbf{Total JEs} &
      \textbf{Cum.\ coverage}\\
    \midrule
    1 & $2$--$9$        & $338$M          & $\phantom{0}92.60\%$\\
    2 & $10$--$99$      & $\phantom{00}23$M & $\phantom{0}99.22\%$\\
    3 & $100$--$999$    & $\phantom{000}2.80$M & $\phantom{0}99.98\%$\\
    4 & $1\,000$--$9\,999$ & $\phantom{000}63$k & $100.00\%$\\
    5 & $\geq 10\,000$  & $\phantom{0000}196$ & $100.00\%$\\
    \bottomrule
  \end{tabular}
\end{table}

\begin{table}[H]
  \centering
  \caption{Detailed distribution within bin~1 ($2$--$9$ line
    items).  Source: \citet{ivertowski2024apn}, Table~III.}
  \label{tab:emp:bin1}
  \begin{tabular}{rrrr}
    \toprule
    \textbf{Line items} & \textbf{Total JEs} & \textbf{Coverage} &
      \textbf{Cum.\ coverage}\\
    \midrule
    $2$ & $221$M          & $60.68\%$ & $60.68\%$\\
    $3$ & $\phantom{00}21$M & $\phantom{0}5.77\%$ & $66.46\%$\\
    $4$ & $\phantom{00}61$M & $16.63\%$ & $83.09\%$\\
    $5$ & $\phantom{00}11$M & $\phantom{0}3.06\%$ & $86.15\%$\\
    $6$ & $\phantom{00}12$M & $\phantom{0}3.32\%$ & $89.47\%$\\
    $7$ & $\phantom{000}4$M & $\phantom{0}1.13\%$ & $90.60\%$\\
    $8$ & $\phantom{000}7$M & $\phantom{0}1.88\%$ & $92.48\%$\\
    $9$ & $\phantom{000}2$M & $\phantom{0}0.42\%$ & $92.90\%$\\
    \bottomrule
  \end{tabular}
\end{table}

\subsection{Debit-credit symmetry}
\label{sec:empirical:symmetry}

Within the corpus, $82\%$ of journal entries have the same number
of debit and credit line items; $11\%$ have fewer credit than debit
items and $7\%$ vice versa.  The dominance of even line counts
follows directly from this symmetry, given that GAAP requires every
posting to balance ($\sum D = \sum C$); the residual $12\%$ with an
odd line count corresponds to entries in which a single side
(debit or credit) carries an aggregated counterpart of multiple
line items on the other.  These statistics motivate Method~A as the
default network-construction strategy in
Section~\ref{sec:network}: any two-line balanced entry admits a
unique $1$-to-$1$ debit-credit edge, and any
$n$-debit-and-$n$-credit balanced entry can be solved without
decomposition.

\subsection{Account population}
\label{sec:empirical:accounts}

The number of distinct general-ledger accounts populated in a single
fiscal year ranges from $76$ to $2\,536$ across the corpus, with
mean $557$ and median $367$
(\citealt{ivertowski2024apn}, Fig.~3).  This range informs
\system{}'s industry-pack expansion (manufacturing, retail,
financial-services, healthcare, technology), which adds sub-account
splits to base charts of accounts to bring synthetic chart sizes into
the empirically observed range.

\subsection{Implications for the benchmark}
\label{sec:empirical:implications}

Three implications drive the design of \bench{}.

\textbf{(i)~Two-line entries dominate.}  Method~A, the default
$1$-to-$1$ debit-credit edge construction, applies to roughly
$60\%$ of all postings and is the single most important
construction step.

\textbf{(ii)~Multi-line entries are the long tail that matters
analytically.}  Although bins 2--5 collectively cover only $7.4\%$
of journal entries, they include the entries that are most
frequently of audit interest---multi-line allocation entries,
period-end accruals, and complex consolidation postings.  Methods
B--E from Section~\ref{sec:network} exist precisely to handle these
cases.

\textbf{(iii)~Even line counts are the norm; odd line counts are
diagnostic.}  The empirical $88\%$/$12\%$ even/odd split, when
combined with audit-relevant flags such as \texttt{is\_post\_close}
and \texttt{document\_type}, supplies a strong baseline against
which structural anomalies can be measured.

The configuration that produces the released \bench{} dataset
(Appendix~\ref{sec:appendix:config}) reproduces these distributions
within the tolerances reported in Section~\ref{sec:evaluation}.
